\section{依法不负刑事责任的精神病人的强制医疗程序}
\begin{enumerate}
  \item 出于避免社会危害和保障精神病患者自身的健康利益而采取的一项对精神病人人身自由予以一定限制并对他的精神疾病进行强制治疗的一种特殊的\zd{安保处分}措施。
  \item 特点。
        \begin{enumerate}
          \item 适用对象特殊。
          \item 应当\zd{经过鉴定}。
          \item 措施带有\zd{强制性}。
          \item 双重保护目的。
        \end{enumerate}
  \item 适用条件。\bs{暴力、无责、有危险}。一定要背出来。
        \begin{enumerate}
          \item 暴力行为严重危害他人人身安全。\al{暴力摧毁垃圾桶 \cuo}。
          \item 应当经过鉴定，不承担刑事责任的精神病人。
          \item 有\bs{继续危害}社会的可能。
        \end{enumerate}
  \item 公安。\bs{跑腿、执行、临时保护措施}。
        \begin{enumerate}
          \item 侦查过程中 \yjt 撤案 \yjt 强疗意见书 \yjt 移送检察院。
          \item 执行法院的强疗\zd{决定}。
          \item \bs{可采取临时性保护约束措施（仅公安）}。
        \end{enumerate}
  \item 检察院。\bs{一审、两不服}。
        \begin{enumerate}
          \item 公安过来的强疗意见书。\bs{先审、再申；问责公安}。
                \begin{enumerate}
                  \item \zd{捕诉部门} 30 日内审查。
                  \item 证据不够要求公安补充 \hz 自行调查。
                  \item 符合。向被申请人居住地、暴力行为所在地的\zd{基层}法院申请强疗。
                  \item 不符合。向公安说明理由。
                  \item 公安没有启动强疗程序。先要求说理，后通知启动。同立案监督。
                \end{enumerate}
          \item 对依申请的强疗决定不服。\bs{20 上法提意见}。20 日内向\zd{上一级法院}提\zd{检察意见}。
          \item 对公诉中法院的强疗判决不服。\bs{判决不服可上诉，检察监督提抗诉}。
        \end{enumerate}
  \item 法院。\bs{自决、一审、一判、一送、四应当}。
        \begin{enumerate}
          \item 审理中自己发现的，可以主动做出决定。
          \item 检察院申请的强疗。\bs{申请强疗 1 月决}。1 个月内作出决定。\yi 同意 \yjt 决定强疗。\er 驳回 \yjt 退回。
          \item 公诉程序中的强疗。\bs{提起公诉要判决}。
                \begin{enumerate}
                  \item 合议庭；原则开庭审理。法定代理人申请 + 法院同意可以不开庭。
                  \item 无责 + 需要强疗 \yjt 判决不负刑事责任 + 强疗决定。
                  \item 无责 + 无需强疗 \yjt 判决不负刑事责任。
                  \item 其他 \yjt 正常处理。
                \end{enumerate}
          \item 权利保障条款的\zd{4 个应当}。背下来。\bs{被申请人、被害人、各自法定代理人、有诉代}。
                \begin{enumerate}
                  \item \zd{通知}被申请人的法定代理人 \hz 其他合适的近亲属到场。
                  \item \zd{会见}被申请人。
                  \item \zd{听取}被害人和法定代理人意见。
                  \item \zd{要有}诉讼代理人，否则就要强制法律援助。
                \end{enumerate}
          \item 决定强疗。\zd{5 日内}将文书送给公安。
        \end{enumerate}
  \item 当事人。\bs{申请复议、申请解除}。
        \begin{enumerate}
          \item \bs{决定不服可复议}。\zd{被害人、法定代理人、近亲属}。\al{上诉 \cuo 仅判决、裁定可上诉}。
          \item 申请解除。
                \begin{enumerate}
                  \item 找谁。谁决定进去的谁决定出来。
                  \item \bs{一个疗程 6 个月}。申请被驳回后，再次申请，间隔 6 个月。
                \end{enumerate}
        \end{enumerate}
  \item 强疗机构。\yi \bs{提解除建议}。\er 解除时，\bs{诊断评估}是否有危险性。
\end{enumerate}
